\documentclass[11pt]{article}
\usepackage[hmargin=2.5cm,vmargin=2.5cm]{geometry}
\usepackage{amsmath,amssymb}
\usepackage{xcolor}
\usepackage[colorlinks=true,linkcolor=blue,urlcolor=blue]{hyperref}

% Light-blue intuition boxes: answers to specific questions raised while
% reviewing the model. Usage: \intuition{Question?}{Answer.}
\definecolor{intuitionblue}{rgb}{0.88,0.94,1.00}
\newcommand{\intuition}[2]{%
  \par\medskip\noindent\hspace*{\fill}%
  \colorbox{intuitionblue}{\parbox{0.94\linewidth}{\textbf{#1}\par\smallskip #2}}%
  \hspace*{\fill}\par\medskip}

% Boxed display for model-code (Dynare) lines: framed box on the current line
% width (so it also fits inside lists), small typewriter font, one equation per
% line (break lines with \\).
\newcommand{\modelcode}[1]{%
  \par\smallskip\noindent\hspace*{\fill}%
  \fbox{\parbox{0.94\linewidth}{\small\ttfamily\raggedright #1}}%
  \hspace*{\fill}\par\smallskip}

\title{Model Derivations:\\ Households, Cost Minimization, Pricing, and Technology}
\author{Spending-efficiency model --- internal note}
\date{\today}

\begin{document}
\maketitle

\noindent
This note derives the model equation by equation, in the order of Section~2 of
the working paper (\texttt{../../draftPaper.tex}): the household's problem, the
intermediate firm's cost minimization, the retailer's Calvo price setting, and
the technology creation-and-adoption block, followed by the government,
market-clearing, balanced-growth, and auxiliary blocks, which contain no
optimization and are stated with brief commentary. The note is complete:
\emph{every} equation of \texttt{models/model\_block.modpart} appears in a
framed box exactly once, at the point where it is derived or stated.
Notation is the paper's throughout. Derivations are in levels; where the
distinction matters, level (un-detrended) variables carry a superscript $\ell$
and plain symbols are the detrended objects the model solves for. The general
detrending logic is in \texttt{../detrending/detrending.tex}, and the technology
trend $1+\gamma^{a}$ is derived in
\texttt{../tech-growth-rate/tech\_growth.tex}.

\tableofcontents

%========================================================================
\section{The Household's Problem}\label{sec:hh}
%========================================================================

\subsection{The Problem}

There is a representative household that consumes, supplies labor, spends time
building human capital, and accumulates private capital and government bonds. Its
problem is:
\begin{equation*}
    \max_{C_t,\,L_t,\,E_t,\,K_t,\,B_t,\,H_t} \;
    \mathbb{E}_0 \sum_{t=0}^{\infty} \beta^{t}
    \left[\ln C_t - \omega\,\frac{\left(L_t + E_t\right)^{1+\varphi}}{1+\varphi}\right]
\end{equation*}
subject to the flow budget constraint,
\begin{equation}\label{eq:budget}
    \left(1+\tau_t^{c}\right) C_t + I_t + B_t
    = \left(1-\tau_t^{w}\right) w_t\, N_t + r_t^{k} K_{t-1}
    + \frac{R_{t-1}}{\Pi_t}\, B_{t-1} + T_t + \Omega_t,
\end{equation}
the private-capital law of motion $K_t = (1-\delta)K_{t-1} + I_t$, the
human-capital law of motion,
\begin{equation}\label{eq:hlom}
    H_t = \left(1-\delta^{h}\right) H_{t-1}
    + \chi\, E_t^{\gamma} \left(\frac{K_{t-1}^{GE}}{\Lambda_{t-1}}\right)^{\mu},
\end{equation}
and the definition of effective labor, $N_t = H_{t-1} L_t$. The household takes
prices, taxes, transfers, and the rebated profits $\Omega_t$ as given. In the
model, the capital law of motion and effective labor read (the $g$ on the
current stock is the stationarization; the human-capital law of motion appears
in Subsection~\ref{subsec:hhg}):
\modelcode{Kp*g = (1-delta)*Kp(-1)+Ip;\\
N = L*H(-1);}
Note two
things about the time margins. First, labor $L_t$ and education time $E_t$ enter
the disutility term jointly, so they compete for the same time endowment. Second,
the human-capital stock enters effective labor with a lag: an hour of schooling
today raises earnings capacity only from tomorrow on.

Substitute investment $I_t = K_t - (1-\delta)K_{t-1}$ and effective labor
$N_t = H_{t-1}L_t$ out of \eqref{eq:budget}, and attach the multiplier $\lambda_t$
to the budget constraint and $\lambda_t^{H}$ to \eqref{eq:hlom}. A Lagrangian for
the household is:
\begin{equation*}
\begin{aligned}
    \mathbb{L} = \mathbb{E}_0 \sum_{t=0}^{\infty} \beta^{t} \Big\{\,
    & \ln C_t - \omega\,\frac{\left(L_t + E_t\right)^{1+\varphi}}{1+\varphi} \\
    & + \lambda_t \Big[\left(1-\tau_t^{w}\right) w_t H_{t-1} L_t
      + r_t^{k} K_{t-1} + \tfrac{R_{t-1}}{\Pi_t} B_{t-1} + T_t + \Omega_t \\
    & \qquad\;\; - \left(1+\tau_t^{c}\right) C_t
      - \big(K_t - (1-\delta) K_{t-1}\big) - B_t \Big] \\
    & + \lambda_t^{H} \Big[\left(1-\delta^{h}\right) H_{t-1}
      + \chi\, E_t^{\gamma} \big(K_{t-1}^{GE}/\Lambda_{t-1}\big)^{\mu}
      - H_t \Big] \Big\}.
\end{aligned}
\end{equation*}

\subsection{First-Order Conditions}

The FOC are:
\begin{align}
    \frac{\partial\mathbb{L}}{\partial C_t} = 0
    &\;\Leftrightarrow\;
    \frac{1}{C_t} = \lambda_t \left(1+\tau_t^{c}\right), \label{eq:focC}\\
    \frac{\partial\mathbb{L}}{\partial B_t} = 0
    &\;\Leftrightarrow\;
    \lambda_t = \beta\,\mathbb{E}_t\!\left[\lambda_{t+1}\,
    \frac{R_t}{\Pi_{t+1}}\right], \label{eq:focB}\\
    \frac{\partial\mathbb{L}}{\partial K_t} = 0
    &\;\Leftrightarrow\;
    \lambda_t = \beta\,\mathbb{E}_t\!\left[\lambda_{t+1}
    \left(r_{t+1}^{k} + 1 - \delta\right)\right], \label{eq:focK}\\
    \frac{\partial\mathbb{L}}{\partial L_t} = 0
    &\;\Leftrightarrow\;
    \omega\left(L_t + E_t\right)^{\varphi}
    = \lambda_t \left(1-\tau_t^{w}\right) w_t H_{t-1}, \label{eq:focL}\\
    \frac{\partial\mathbb{L}}{\partial E_t} = 0
    &\;\Leftrightarrow\;
    \omega\left(L_t + E_t\right)^{\varphi}
    = \lambda_t^{H}\, \chi\gamma\, E_t^{\gamma-1}
    \big(K_{t-1}^{GE}/\Lambda_{t-1}\big)^{\mu}, \label{eq:focE}\\
    \frac{\partial\mathbb{L}}{\partial H_t} = 0
    &\;\Leftrightarrow\;
    \lambda_t^{H} = \beta\,\mathbb{E}_t\!\left[\lambda_{t+1}
    \left(1-\tau_{t+1}^{w}\right) w_{t+1} L_{t+1}
    + \lambda_{t+1}^{H}\left(1-\delta^{h}\right)\right]. \label{eq:focH}
\end{align}
For \eqref{eq:focK}, note that $K_t$ appears twice in the Lagrangian: bought today
(through the substituted investment, at utility price $\lambda_t$) and carried
into tomorrow's income and resale (at $\lambda_{t+1}(r^k_{t+1}+1-\delta)$). For
\eqref{eq:focH} the structure is the same: an extra unit of $H_t$ raises
tomorrow's effective labor --- and hence after-tax labor income
$\lambda_{t+1}(1-\tau^w_{t+1})w_{t+1}L_{t+1}$ --- and survives into the future
stock at rate $1-\delta^{h}$, valued at $\lambda^{H}_{t+1}$.

\eqref{eq:focC} identifies $\lambda_t$ as the marginal utility of wealth.
\eqref{eq:focB}--\eqref{eq:focK} are the bond and capital Euler equations; the
$g_{t+1}$ in their model counterparts below is the stationarization (see
Subsection~\ref{subsec:hhg}). \eqref{eq:focL} is the labor-supply condition: an
hour of work supplies $H_{t-1}$ efficiency units paid $w_t$ each, so its utility
value is the after-tax hourly wage times the marginal utility of wealth. In the
model (\texttt{models/model\_block.modpart}), the six conditions read:
\modelcode{1/C = lambda*(1+tauc);\\
lambda = betta*(lambda(+1)/g(+1)*R/PI(+1));\\
1 = betta*(lambda(+1)/lambda/g(+1)*(1-delta+rk(+1)));\\
omega*(L+E)\^{}varphi = lambda*w*H(-1)*(1-tauw);\\
omega*(L+E)\^{}varphi = lambda\_H*chiH*gamma*E\^{}(gamma-1)*(Kge(-1))\^{}(mu);\\
lambda\_H = betta*(lambda(+1)*(1-tauw(+1))*w(+1)*L(+1)+lambda\_H(+1)*(1-deltaH));}

\subsection{The Work-vs-Schooling Arbitrage}

Because $L_t$ and $E_t$ draw on the same time endowment, their marginal
disutilities coincide --- the left-hand sides of \eqref{eq:focL} and
\eqref{eq:focE} are identical. Equating the right-hand sides:
\begin{equation}\label{eq:arbitrage}
    \underbrace{\lambda_t \left(1-\tau_t^{w}\right) w_t
    H_{t-1}}_{\text{return to an hour of work}}
    \;=\;
    \underbrace{\lambda_t^{H}\, \chi\gamma\, E_t^{\gamma-1}
    \big(K_{t-1}^{GE}/\Lambda_{t-1}\big)^{\mu}}_{\text{return to an hour of
    schooling}}.
\end{equation}
\eqref{eq:arbitrage} says the household allocates time until the after-tax wage
of the marginal hour equals the shadow value of the human capital that hour of
schooling builds. This is the margin the public human-capital experiments work
through: a higher $K^{GE}$ raises the marginal product of schooling time, tilting
\eqref{eq:arbitrage} toward education, and the resulting $H$ raises effective
labor with a lag.

\subsection{The Stochastic Discount Factor}\label{subsec:sdf}

The household's valuation of a real payoff at $t+1$ per unit at $t$ is the ratio
of marginal utilities of \emph{income}. From \eqref{eq:focC}, a unit of income at
$t$ buys $1/(1+\tau^c_t)$ units of consumption, so:
\begin{equation}\label{eq:sdf}
    \mathit{SDF}_{t+1}
    = \beta\,\frac{\lambda_{t+1}\left(1+\tau_{t+1}^{c}\right)}
    {\lambda_t\left(1+\tau_t^{c}\right)}
    = \beta\,\frac{C_t}{C_{t+1}},
\end{equation}
where the second equality substitutes \eqref{eq:focC}; it is exact, with no
constant-tax assumption needed. In the model:
\modelcode{SDF = betta*lambda*(1+tauc)/(lambda(-1)*(1+tauc(-1)));}
The SDF is
defined by household preferences but used elsewhere: by the technology adopters
(Section~\ref{sec:tech}) and, in the $\beta^{s}\lambda_{t+s}$ form, by the
price-setting retailers (Section~\ref{sec:pricing}).

\subsection{Aside: Which Conditions Pick Up $g$ When Stationarized}\label{subsec:hhg}

Detrending maps $\lambda_t \mapsto \hat\lambda_t/\Lambda_t$ (marginal utility
falls as the economy grows) and leaves $H_t$, $L_t$, and $E_t$ untouched (all
stationary). The consequences, condition by condition:
\begin{itemize}
    \item \textbf{Bond and capital Eulers \eqref{eq:focB}--\eqref{eq:focK}: a
    $1/g_{t+1}$ appears.} The two $\lambda$'s sit at different dates, so their
    trends do not cancel: $\lambda_{t+1}/\lambda_t =
    (\hat\lambda_{t+1}/\hat\lambda_t)/g_{t+1}$.
    \item \textbf{Labor FOC \eqref{eq:focL}: nothing.} $\lambda_t w_t$ is a
    same-date product of an inverse-trend and a trend variable --- stationary.
    \item \textbf{Education and human-capital FOCs
    \eqref{eq:focE}--\eqref{eq:focH}: nothing.} $H$ is a stationary stock, so its
    shadow value $\lambda^{H}$ is stationary too ($\lambda_{t+1}w_{t+1}L_{t+1}$
    is again an inverse-trend times a trend). This is why the model's
    \texttt{lambda\_H} equation carries no $g$ while the capital Euler does.
    \item \textbf{The HC technology in \eqref{eq:hlom}: the $\Lambda_{t-1}$
    scaling.} Public capital $K^{GE}$ trends while $H$ does not; dividing by the
    trend inside the accumulation technology is what keeps \eqref{eq:hlom}
    balanced. \texttt{Kge} is already the detrended stock, so in the model the
    equation reads:
    \modelcode{H = (1-deltaH)*H(-1)+chiH*E\^{}gamma*(Kge(-1))\^{}(mu);}
\end{itemize}

%========================================================================
\section{Intermediate Goods: Cost Minimization and the Markup Wedge}\label{sec:costmin}
%========================================================================

This section derives the factor-demand conditions of Appendix~A in the working
paper from the intermediate firm's cost-minimization problem, and shows where the
fixed gross markup $\mu^{p}$ (model parameter \texttt{markupss})
enters.\footnote{An earlier version of this derivation carried a fixed cost of
production $\Theta$; the parameter was calibrated to zero and has since been
removed from the model and the paper entirely.}

\subsection{The Firm's Problem}

There is a continuum of intermediate producers indexed by $i$. A typical firm
produces output using private capital and effective labor, and takes as given the
factor prices $r_t^{k}$ and $w_t$ as well as two externalities: the
public-infrastructure stock $K_{t-1}^{GI}$ and the adopted-technology stock
$A_{t-1}$. Its technology is
\begin{equation}\label{eq:tech}
    Y_t(i) = A_{t-1}^{\vartheta-1}\left(K_{t-1}^{GI}\right)^{\alpha_G}
    K_{t-1}(i)^{\alpha}\,N_t(i)^{1-\alpha}.
\end{equation}
In the model, the aggregate production function (tagged \texttt{[name='y']})
reads:
\modelcode{y = A(-1)\^{}(vartheta-1)*(Kg(-1)\^{}(alphaG))*(Kp(-1)\^{}alpha)*(N\^{}(1-alpha));}

Consider the problem of producing a required level of output, $\bar{Y}_t(i)$, as
cheaply as possible. The firm's problem is:
\begin{equation}\label{eq:cmproblem}
    \min_{K_{t-1}(i),\,N_t(i)} \; r_t^{k}\,K_{t-1}(i) + w_t\,N_t(i)
\end{equation}
subject to:
\begin{equation}\label{eq:cmconstraint}
    A_{t-1}^{\vartheta-1}\left(K_{t-1}^{GI}\right)^{\alpha_G}
    K_{t-1}(i)^{\alpha}\,N_t(i)^{1-\alpha} \;\ge\; \bar{Y}_t(i).
\end{equation}
Note that the required output level is what makes the problem well-posed: if we
dropped it, the cost-minimizing input choice would just be zero.

\subsection{Lagrangian and First-Order Conditions}

A Lagrangian for the firm is:
\begin{equation*}
    \mathbb{L} = -\,r_t^{k}\,K_{t-1}(i) - w_t\,N_t(i)
    + \xi_t\left[A_{t-1}^{\vartheta-1}\left(K_{t-1}^{GI}\right)^{\alpha_G}
      K_{t-1}(i)^{\alpha}\,N_t(i)^{1-\alpha} - \bar{Y}_t(i)\right].
\end{equation*}
The FOC are:
\begin{align*}
    \frac{\partial\mathbb{L}}{\partial K_{t-1}(i)} = 0
    &\;\Leftrightarrow\; r_t^{k} = \xi_t\,\alpha\,\frac{Y_t(i)}{K_{t-1}(i)},\\
    \frac{\partial\mathbb{L}}{\partial N_t(i)} = 0
    &\;\Leftrightarrow\; w_t = \xi_t\,(1-\alpha)\,\frac{Y_t(i)}{N_t(i)}.
\end{align*}
Dividing the first condition by the second, the multiplier and the output term drop
out:
\begin{equation}\label{eq:mix}
    \frac{K_{t-1}(i)}{N_t(i)} = \frac{\alpha}{1-\alpha}\,\frac{w_t}{r_t^{k}}.
\end{equation}
\eqref{eq:mix} says that the cost-minimizing capital-labor ratio depends only on
factor prices. Nothing on the right-hand side depends on $i$, so all firms choose
the same mix regardless of their scale. In the model:
\modelcode{Kp(-1)/N = alpha/(1-alpha)*w/rk;}

\subsection{The Multiplier Is Marginal Cost; the Markup Wedge}

What is $\xi_t$? It is the multiplier on the output constraint: it measures how
much the minimized cost rises if the firm must produce one more unit. That is just
the firm's \emph{real marginal cost}. Since the technology has constant returns in
the two chosen inputs, marginal cost equals unit factor cost.

The firm sells its output to the retailers at a fixed gross markup $\mu^{p}$ over
marginal cost. The relative price it receives is $mc_t$ --- which is in turn the
retailers' marginal cost --- so:
\begin{equation}\label{eq:wedge}
    mc_t = \mu^{p}\,\xi_t
    \qquad\Leftrightarrow\qquad
    \xi_t = \frac{mc_t}{\mu^{p}}.
\end{equation}
Plugging \eqref{eq:wedge} into the FOC, we get the factor-demand conditions:
\begin{align}
    r_t^{k} &= \alpha\,\frac{mc_t}{\mu^{p}}\,\frac{Y_t}{K_{t-1}},\label{eq:fdK}\\
    w_t &= (1-\alpha)\,\frac{mc_t}{\mu^{p}}\,\frac{Y_t}{N_t}.\label{eq:fdN}
\end{align}
\eqref{eq:fdK}--\eqref{eq:fdN} are the factor demands in Appendix~A of the paper.
\eqref{eq:fdN} is the model's labor-demand equation, and the steady state
(\texttt{models/Steady\_states\_solution.m}) solves it the same way:
\modelcode{(1-alpha)*mc*y/N = markupss*w;\\
w=(1-alpha)*mc*y/N/markupss;}

Now consider profits. Per unit of output the firm receives $mc_t$ and pays
$mc_t/\mu^{p}$ in factor costs, so the profit flow is:
\begin{equation}\label{eq:profit}
    \frac{\mu^{p}-1}{\mu^{p}}\,mc_t\,Y_t.
\end{equation}
\eqref{eq:profit} is the dividend that accrues to the owners of adopted
technologies --- the first term in the value function $\mathcal{V}_t$
(Section~\ref{sec:tech}). This is why the wedge in
\eqref{eq:fdK}--\eqref{eq:fdN} matters: with competitive factor pricing at $mc_t$,
factor payments would exhaust revenue and there would be no profit flow left to
reward technology adoption.

\subsection{Aside: Where Is the Pricing Equation?}

Where in the model can one see that the firm sells at the markup $\mu^{p}$ over
marginal cost? Nowhere as a standalone line: the model never writes the
intermediate firm's pricing rule as its own equation. Instead, \eqref{eq:wedge} is
embedded in the labor-demand equation. Take the model's
\texttt{(1-alpha)*mc*y/N = markupss*w} and solve it for $mc_t$:
\begin{equation}\label{eq:pricingrule}
    mc_t = \mu^{p}\,\frac{w_t}{(1-\alpha)\,Y_t/N_t}.
\end{equation}
The second factor is the wage divided by the marginal product of labor --- the cost
of producing one more unit using labor, which is just marginal cost $\xi_t$. So
\eqref{eq:pricingrule} is the same equation as \eqref{eq:fdN}, read the other way
around: as a labor demand it prices labor off $mc_t/\mu^{p}$; as a pricing rule it
sets the output price at $\mu^{p}$ times marginal cost. One equation carries both
statements, which is why there is no separate pricing line to point to.

Two other equations pin down the interpretation. First, that $mc_t$ is the price
the retailers pay is visible in the Calvo recursion
\texttt{x1 = lambda*mc*yd + ...}: $mc_t$ enters as the \emph{retailers'} real
marginal cost, and since retailers transform intermediate output one-for-one,
their marginal cost is exactly the relative price at which they buy it. Second,
the dividend in the technology value, \eqref{eq:profit}, equals
$(mc_t - mc_t/\mu^{p})\,Y_t$ --- price minus unit cost times quantity, exactly the
monopoly profit of a producer charging the fixed gross markup $\mu^{p}$. If the
wedge in the factor demands were something else (a tax, say), this profit would
not flow to the technology owners. The three equations together --- labor demand,
Calvo recursion, technology value --- are what make the fixed-markup reading the
right one.

\subsection{Aside: Two Markups, Two Layers}

$\mu^{p}$ (\texttt{markupss} $=1.18$) is a calibrated parameter and is \emph{not}
the retail Calvo markup $\epsilon/(\epsilon-1) = 10/9 \approx 1.11$. The two
markups belong to different layers of the production side. $\mu^{p}$ is charged by
the intermediate producers, and their profits reward technology adoption.
$\epsilon$ governs the retailers, whose staggered price setting generates the New
Keynesian block. Keeping them separate is what lets the model calibrate the profit
share of the technology block independently of the nominal rigidity.

%========================================================================
\section{Price Setting}\label{sec:pricing}
%========================================================================

This section derives the four pricing equations of the model --- the two auxiliary
recursions for $x_{1,t}$ and $x_{2,t}$, the reset-price condition
$\epsilon x_{1,t} = (\epsilon-1)x_{2,t}$, and the price-index law of motion ---
from the retailer's problem.

\subsection{Setup}

Retailers purchase intermediate output at the relative price $mc_t$ --- their real
marginal cost (Section~\ref{sec:costmin}) --- differentiate it one-for-one into
varieties, and sell variety $i$ at the nominal price $P_t(i)$. The final-good
producer's cost minimization gives the demand for each variety:
\begin{equation}\label{eq:demand}
    y_t(i) = \left(\frac{P_t(i)}{P_t}\right)^{-\epsilon} Y_t^{d}.
\end{equation}
Prices are sticky \`a la Calvo: each period a retailer reoptimizes its price with
probability $1-\theta_p$. With probability $\theta_p$ it cannot, and instead
indexes its price to past inflation with elasticity $\chi_p$, so a non-reoptimized
price is scaled by $\Pi_{t-1}^{\chi_p}$ each period.

Consider a retailer that reoptimizes at $t$. If it is still stuck with that price
at $t+s$, indexation will have scaled it by the cumulative factor
\begin{equation}\label{eq:index}
    X_{t,s} = \prod_{k=1}^{s} \Pi_{t+k-1}^{\chi_p}, \qquad X_{t,0} = 1.
\end{equation}

\subsection{The Reset-Price Problem}

The retailer chooses $P_t(i)$ to maximize the expected discounted real profits it
earns while the price remains in force. A profit $s$ periods ahead survives with
probability $\theta_p^{s}$ and is valued with the household's discount factor
$\beta^{s}\lambda_{t+s}/\lambda_t$; multiplying through by the constant
$\lambda_t$, the problem can be written:
\begin{equation*}
    \max_{P_t(i)} \; \mathbb{E}_t \sum_{s=0}^{\infty}
    \left(\beta\theta_p\right)^{s} \lambda_{t+s}
    \left[\frac{P_t(i)\,X_{t,s}}{P_{t+s}} - mc_{t+s}\right]
    \left(\frac{P_t(i)\,X_{t,s}}{P_{t+s}}\right)^{-\epsilon} Y_{t+s}^{d}.
\end{equation*}
This is why $\lambda$ appears inside the model's auxiliaries below: profits are
valued in utility units. The FOC with respect to $P_t(i)$ is:
\begin{equation*}
    \mathbb{E}_t \sum_{s=0}^{\infty} \left(\beta\theta_p\right)^{s} \lambda_{t+s}
    \left[(1-\epsilon)\,P_t(i)^{-\epsilon}
    \left(\frac{X_{t,s}}{P_{t+s}}\right)^{1-\epsilon}
    + \epsilon\, mc_{t+s}\,P_t(i)^{-\epsilon-1}
    \left(\frac{X_{t,s}}{P_{t+s}}\right)^{-\epsilon}\right] Y_{t+s}^{d} = 0.
\end{equation*}
Multiplying by $P_t(i)^{\epsilon+1}$ and rearranging:
\begin{equation}\label{eq:focP}
    (\epsilon-1)\,P_t(i)\;\mathbb{E}_t \sum_{s=0}^{\infty}
    \left(\beta\theta_p\right)^{s} \lambda_{t+s}
    \left(\frac{X_{t,s}}{P_{t+s}}\right)^{1-\epsilon} Y_{t+s}^{d}
    = \epsilon\;\mathbb{E}_t \sum_{s=0}^{\infty}
    \left(\beta\theta_p\right)^{s} \lambda_{t+s}\, mc_{t+s}
    \left(\frac{X_{t,s}}{P_{t+s}}\right)^{-\epsilon} Y_{t+s}^{d}.
\end{equation}
Nothing in \eqref{eq:focP} depends on $i$, so all reoptimizing retailers choose the
same reset price; call it $P_t^{\#}$.

\subsection{The Recursive Form}

\eqref{eq:focP} contains infinite sums and the price level; to obtain the model's
stationary system, express everything in relative terms. Define the relative reset
price and the indexation-deflated price drift:
\begin{equation}\label{eq:defs}
    \Pi_t^{*} \equiv \frac{P_t^{\#}}{P_t}, \qquad
    D_{t,s} \equiv \frac{X_{t,s}\,P_t}{P_{t+s}}, \qquad D_{t,0} = 1.
\end{equation}
$D_{t,s}$ is the relative price at $t+s$ of a retailer that reset to the aggregate
price level at $t$ and indexed ever since. Splitting off the first indexation step,
it obeys:
\begin{equation}\label{eq:Drec}
    D_{t,s} = \frac{\Pi_t^{\chi_p}}{\Pi_{t+1}}\; D_{t+1,s-1}.
\end{equation}
Now define the two sums in \eqref{eq:focP} in relative terms:
\begin{equation}\label{eq:x1def}
    x_{1,t} \equiv \mathbb{E}_t \sum_{s=0}^{\infty}
    \left(\beta\theta_p\right)^{s} \lambda_{t+s}\, mc_{t+s}\, D_{t,s}^{-\epsilon}\,
    Y_{t+s}^{d}, \qquad
    f_{t} \equiv \mathbb{E}_t \sum_{s=0}^{\infty}
    \left(\beta\theta_p\right)^{s} \lambda_{t+s}\, D_{t,s}^{1-\epsilon}\,
    Y_{t+s}^{d}.
\end{equation}
$x_{1,t}$ discounts marginal cost and $f_t$ discounts revenue, both weighted by the
demand the price will face. Using \eqref{eq:Drec}, splitting off the $s=0$ term
turns each sum into a recursion:
\begin{equation}\label{eq:x1}
    x_{1,t} = \lambda_t\, mc_t\, Y_t^{d}
    + \beta\theta_p\, \mathbb{E}_t\!\left[
    \left(\frac{\Pi_t^{\chi_p}}{\Pi_{t+1}}\right)^{-\epsilon} x_{1,t+1}\right],
\end{equation}
\begin{equation}\label{eq:f}
    f_{t} = \lambda_t\, Y_t^{d}
    + \beta\theta_p\, \mathbb{E}_t\!\left[
    \left(\frac{\Pi_t^{\chi_p}}{\Pi_{t+1}}\right)^{1-\epsilon} f_{t+1}\right].
\end{equation}
The model carries the revenue side as $x_{2,t} \equiv \Pi_t^{*} f_t$. Multiplying
\eqref{eq:f} by $\Pi_t^{*}$:
\begin{equation}\label{eq:x2}
    x_{2,t} = \lambda_t\, \Pi_t^{*}\, Y_t^{d}
    + \beta\theta_p\, \mathbb{E}_t\!\left[
    \left(\frac{\Pi_t^{\chi_p}}{\Pi_{t+1}}\right)^{1-\epsilon}
    \frac{\Pi_t^{*}}{\Pi_{t+1}^{*}}\, x_{2,t+1}\right].
\end{equation}
Finally, dividing \eqref{eq:focP} by $P_t$ and substituting the definitions, the
FOC becomes:
\begin{equation}\label{eq:reset}
    \epsilon\, x_{1,t} = (\epsilon-1)\, x_{2,t}.
\end{equation}
\eqref{eq:x1}, \eqref{eq:x2}, and \eqref{eq:reset} are the model equations:
\modelcode{x1 = lambda*mc*yd+betta*thetap*(PI\^{}chi/PI(+1))\^{}(-epsilon)*x1(+1);\\
x2 = lambda*PIstar*yd+betta*thetap*(PI\^{}chi/PI(+1))\^{}(1-epsilon)*PIstar/PIstar(+1)*x2(+1);\\
epsilon*x1 = (epsilon-1)*x2;}
\eqref{eq:reset} says the reset price is a
markup over (discounted) marginal cost: absent stickiness
($\theta_p=0$) it collapses to
$\Pi_t^{*} = \tfrac{\epsilon}{\epsilon-1}\,mc_t$, the static markup rule.

\subsection{The Price-Index Law of Motion}

The price index aggregates $P_t^{1-\epsilon} = \int_0^1 P_t(i)^{1-\epsilon}di$.
Each period a fraction $\theta_p$ of prices are last period's scaled by
$\Pi_{t-1}^{\chi_p}$ --- by the Calvo assumption the non-reoptimizers are a random
sample, so their price distribution is last period's --- and a fraction
$1-\theta_p$ equal $P_t^{\#}$:
\begin{equation*}
    P_t^{1-\epsilon}
    = \theta_p \left(\Pi_{t-1}^{\chi_p}\right)^{1-\epsilon} P_{t-1}^{1-\epsilon}
    + \left(1-\theta_p\right) \left(P_t^{\#}\right)^{1-\epsilon}.
\end{equation*}
Dividing by $P_t^{1-\epsilon}$:
\begin{equation}\label{eq:pindex}
    1 = \theta_p \left(\frac{\Pi_{t-1}^{\chi_p}}{\Pi_t}\right)^{1-\epsilon}
    + \left(1-\theta_p\right)\left(\Pi_t^{*}\right)^{1-\epsilon},
\end{equation}
in the model:
\modelcode{1 = thetap*(PI(-1)\^{}chi/PI)\^{}(1-epsilon)+(1-thetap)*PIstar\^{}(1-epsilon);}
\eqref{eq:pindex} determines inflation: given
the reset price from \eqref{eq:reset}, it pins down $\Pi_t$.

\subsection{Aside: Price Dispersion Is an Aggregation Object}\label{subsec:vp}

The same split between reoptimizers and indexers, applied to the demand weights
$\left(P_t(i)/P_t\right)^{-\epsilon}$ instead of $P_t(i)^{1-\epsilon}$, gives the
law of motion of price dispersion $v_t^{p} \equiv \int_0^1
\left(P_t(i)/P_t\right)^{-\epsilon} di$:
\begin{equation}\label{eq:vp}
    v_t^{p} = \theta_p \left(\frac{\Pi_{t-1}^{\chi_p}}{\Pi_t}\right)^{-\epsilon}
    v_{t-1}^{p} + \left(1-\theta_p\right)\left(\Pi_t^{*}\right)^{-\epsilon}.
\end{equation}
In the model:
\modelcode{vp = thetap*(PI(-1)\^{}chi/PI)\^{}(-epsilon)*vp(-1)+(1-thetap)*PIstar\^{}(-epsilon);}
We derive it here because the algebra is one more application of the same split,
but $v_t^{p}$ is \emph{not} part of the pricing block: no retailer chooses
anything with respect to it. It is an aggregation statistic --- summing the variety
demands \eqref{eq:demand} gives $\int_0^1 y_t(i)\,di = v_t^{p}\,Y_t^{d} = Y_t$,
the wedge between gross output and demand --- and accordingly the paper and the
glossary keep it with market clearing.

\subsection{Aside: Why No Growth Factor in the Recursions}\label{subsec:prg}

Unlike the Euler equations, the recursions \eqref{eq:x1} and \eqref{eq:x2} carry no
growth factor in the model. The reason is that $\lambda$ and $Y^{d}$ enter only as
the product $\lambda_{t+s} Y_{t+s}^{d}$: the marginal utility scales inversely
with the trend and demand proportionally to it, so the product is stationary and
the detrending factors cancel term by term.

%========================================================================
\section{Technology Creation and Adoption}\label{sec:tech}
%========================================================================

This section derives the technology block: the creation process for $Z_t$, the
adoption law of motion for $A_t$, the values $\mathcal{V}_t$ and $\mathcal{J}_t$
of adopted and unadopted technologies, and the FOC for adoption expenditure
$S_t$. The dividend that makes technologies valuable is \eqref{eq:profit} from
Section~\ref{sec:costmin}.

\subsection{Setup}

There are two stocks. $Z_t$ is the stock of \emph{created} technologies ---
inventions that exist but are not yet usable in production. $A_t$ is the stock of
\emph{adopted} technologies, which enters production through the love-of-variety
term $A_{t-1}^{\vartheta-1}$. The gap $Z_t - A_t > 0$ is the pool of unadopted
technologies. Competitive adopters spend resources to convert unadopted
technologies into adopted ones; each adopted technology earns the per-period
profit flow of the intermediate producers, which is what gives technologies their
value. Every technology, adopted or not, survives from one period to the next
with probability $\phi$.

\subsection{Technology Creation}

Creation is a reduced-form process, not an optimization: created technology
responds to lagged human capital (a learning-by-doing channel) and to lagged
\emph{effective} public R\&D spending,
\begin{equation}\label{eq:Z}
    \ln\frac{Z_t}{Z^{SS}} = \rho_A \ln\frac{Z_{t-1}}{Z^{SS}}
    + \alpha_{HA}\,\ln\frac{H_{t-1}}{H^{SS}}
    + \alpha_{RD}\,\ln\frac{\left(1-e^{GRD}_{t-1}\right) G^{RD}_{t-1}}
    {\left(1-e^{GRD}\right) G^{RD,SS}},
\end{equation}
in the model:
\modelcode{ln(Z/STEADY\_STATE(Z)) = rho\_A*ln(Z(-1)/STEADY\_STATE(Z))\\
\hspace*{2em}+alphaRD*ln((1-eGRD(-1))*Grd(-1)/((1-eGRD\_ss)*STEADY\_STATE(Grd)))\\
\hspace*{2em}+alphaHA*ln(H(-1)/STEADY\_STATE(H));}
Three remarks. First,
$\alpha_{HA}$ and $\alpha_{RD}$ are \emph{per-period loadings}; because the
process is an AR(1) with persistence $\rho_A$, the long-run elasticities of
created technology are $\alpha_{HA}/(1-\rho_A)$ and $\alpha_{RD}/(1-\rho_A)$.
Second, only the productive fraction $1-e^{GRD}_{t-1}$ of R\&D spending drives
creation, while the raw spending $G^{RD}_t$ enters the government budget --- the
efficiency wedge destroys the difference. Third, because \eqref{eq:Z} is written
in ratios to steady state, it looks identical in levels and in detrended
variables: the common trend cancels.

\subsection{The Adoption Law of Motion}

Between $t-1$ and $t$, a fraction $\phi$ of all technologies survives, and each
surviving unadopted technology becomes adopted with probability $q_t$. In levels:
\begin{equation}\label{eq:Alevel}
    A^{\ell}_t = q_t\,\phi\left(Z^{\ell}_{t-1} - A^{\ell}_{t-1}\right)
    + \phi\,A^{\ell}_{t-1},
\end{equation}
newly adopted technologies plus the surviving previously adopted stock. The
technology stocks grow at their own gross rate
$1+\gamma^{a} = g^{(1-\alpha)/(\vartheta-1)}$, faster than output. Substituting
$A^{\ell}_t = (1+\gamma^{a})^{t} A_t$ and $Z^{\ell}_t = (1+\gamma^{a})^{t} Z_t$
into \eqref{eq:Alevel} and dividing by $(1+\gamma^{a})^{t-1}$:
\begin{equation}\label{eq:Astat}
    \left(1+\gamma^{a}\right) A_t
    = q_t\,\phi\left(Z_{t-1} - A_{t-1}\right) + \phi\,A_{t-1},
\end{equation}
in the model:
\modelcode{(1+gammaa)*A = q*phi*(Z(-1)-A(-1))+phi*A(-1);}
The
current-dated stock picks up the growth factor, exactly as the capital law of
motion picks up $g_t$.

\intuition{Why does the survival rate multiply the unadopted pool?}{%
Obsolescence is a property of the technology itself, not of its adoption
status: $\phi$ is the probability that an idea remains economically relevant
one period on, and that fate can befall a blueprint on the shelf just as much
as a technology in production. Of the unadopted pool $Z_{t-1}-A_{t-1}$, only
the surviving fraction $\phi$ is still around to be adopted, and each survivor
is adopted with probability $q_t$ --- one cannot adopt a dead idea. Dropping
the $\phi$ would make the shelf an obsolescence-proof vault: leaving a
technology unadopted would shelter it from mortality, artificially raising the
option value of waiting and distorting the adoption margin; with a uniform
$\phi$, delay carries the genuine risk that the idea dies unused, which is
what makes adoption expenditure urgent. Uniform survival also keeps the pool
arithmetic coherent ($A$ is a subset of $Z$, so both must inherit the same
mortality), and it is the same $\phi$ that discounts the continuation values
in \eqref{eq:jsmall}: the $\phi\, q_t$ product here is exactly the event the
adopter pays $S_t$ to bet on --- survive, then convert. The single rate for
adopted and unadopted technologies follows Comin--Gertler (2006) and
Anzoategui et al.\ (2019); \texttt{phi} $=0.98$ quarterly, about 8 percent
obsolescence per year.}

\subsection{The Adoption Probability}

Each unadopted technology is adopted next period with probability
\begin{equation}\label{eq:q}
    q_t = \left(q_0 + \varepsilon^{q}_t\right)\left(S_t\right)^{\varsigma},
\end{equation}
increasing
in adoption expenditure $S_t$ with elasticity $\varsigma$; in the model:
\modelcode{q = (kappaprob+epsi\_q)*(S)\^{}(varsigma);}
Two conventions hide in
\eqref{eq:q}. First, $S_t$ is already detrended, so no $\Lambda_t$ appears (the
paper's Section~2 writes the level form $q_t = q_0(S^{\ell}_t/\Lambda_t)^{\varsigma}$).
Second, as the next subsection makes precise, $S_t$ is expenditure \emph{per
unadopted technology, scaled by the adopted stock} --- the same scaling the value
functions use.

\subsection{The Value of an Adopted Technology}

Start from one adopted variety. It earns the per-variety profit $\pi_t$, and its
owner values the surviving continuation with the household stochastic discount
factor (Subsection~\ref{subsec:sdf}), so in levels:
\begin{equation}\label{eq:vsmall}
    v_t = \pi_t + \phi\,\mathbb{E}_t\!\left[\mathit{SDF}^{\ell}_{t+1}\,
    v_{t+1}\right],
    \qquad
    \pi_t = \frac{1}{A^{\ell}_{t-1}}\,\frac{\mu^{p}-1}{\mu^{p}}\,mc_t\,Y^{\ell}_t,
\end{equation}
where the second expression divides the total intermediate-producer profit flow
\eqref{eq:profit} by the $A^{\ell}_{t-1}$ varieties active at $t$.

\intuition{What is the microfoundation of the value recursion in \eqref{eq:vsmall}?}{%
It is not anyone's optimality condition over a quantity; it is an
\emph{asset-pricing} (no-arbitrage) condition --- the household's portfolio FOC
for claims to variety profits. An adopted technology is a durable, tradeable
claim: the exclusive right to produce variety $i$ and collect its monopoly rent
$\pi_t$, itself microfounded in Section~\ref{sec:costmin} (love-of-variety
demand gives market power; the markup $\mu^{p}$ converts it into the rent).
Households ultimately own these claims --- producer profits are rebated through
$\Omega_t$ in the budget constraint --- so the claim is priced with the
household's SDF (Subsection~\ref{subsec:sdf}). Let $p_t$ be the ex-dividend
price; adding holdings of the claim to the household's budget and taking the
FOC, with the claim paying off only if the technology survives (probability
$\phi$, worth zero otherwise), gives
$p_t = \phi\,\mathbb{E}_t[\mathit{SDF}^{\ell}_{t+1}(\pi_{t+1}+p_{t+1})]$.
Defining the cum-dividend value $v_t \equiv \pi_t + p_t$ and substituting
yields \eqref{eq:vsmall}. The structure is the standard Lucas-tree pricing of a
dividend stream with survival risk, and it is the same machinery the model uses
for its other long-lived objects: the human-capital shadow value
\eqref{eq:focH} and the Calvo reset-price recursions
\eqref{eq:x1}--\eqref{eq:x2}.}

The model does not carry $v_t$; it carries the value \emph{scaled by the adopted
stock}, $\mathcal{V}^{\ell}_t \equiv A^{\ell}_{t-1} v_t$. Multiplying
\eqref{eq:vsmall} by $A^{\ell}_{t-1}$:
\begin{equation}\label{eq:Vlevel}
    \mathcal{V}^{\ell}_t = \frac{\mu^{p}-1}{\mu^{p}}\,mc_t\,Y^{\ell}_t
    + \phi\,\mathbb{E}_t\!\left[\mathit{SDF}^{\ell}_{t+1}\,
    \frac{A^{\ell}_{t-1}}{A^{\ell}_{t}}\,\mathcal{V}^{\ell}_{t+1}\right],
\end{equation}
where the ratio $A^{\ell}_{t-1}/A^{\ell}_{t}$ appears because
$\mathcal{V}^{\ell}_{t+1} = A^{\ell}_{t}\,v_{t+1}$ is scaled by \emph{next}
period's active stock. This is where the $A_{t-1}/A_t$ factor in the model comes
from: it converts between the two scalings, nothing more.

Now detrend. $\mathcal{V}^{\ell}$ and the profit flow grow with the output trend
$\Lambda_t$; the level SDF relates to the model's detrended one by
$\mathit{SDF}^{\ell}_{t+1} = \mathit{SDF}_{t+1}/g_{t+1}$; and the technology-stock
ratio converts as $A^{\ell}_{t-1}/A^{\ell}_t =
(A_{t-1}/A_t)/(1+\gamma^{a})$. Substituting into \eqref{eq:Vlevel} and dividing by
$\Lambda_t$, the $g_{t+1}$ from the SDF cancels against the trend growth
$\Lambda_{t+1}/\Lambda_t$ of the continuation value, leaving:
\begin{equation}\label{eq:V}
    \mathcal{V}_t = \frac{\mu^{p}-1}{\mu^{p}}\,mc_t\,Y_t
    + \phi\,\mathbb{E}_t\!\left[\mathit{SDF}_{t+1}\,
    \frac{A_{t-1}}{A_{t}}\,\frac{1}{1+\gamma^{a}}\,\mathcal{V}_{t+1}\right],
\end{equation}
in the model:
\modelcode{V = (markupss-1)/(markupss)*mc*y + phi*SDF(+1)*V(+1)*A(-1)/A/(1+gammaa);}

\subsection{The Value of an Unadopted Technology}

An unadopted technology costs its adopter $s_t$ this period; conditional on
surviving (probability $\phi$), it becomes adopted next period with probability
$q_t$ --- worth $v_{t+1}$ --- and stays unadopted otherwise --- worth $j_{t+1}$.
In levels, per variety:
\begin{equation}\label{eq:jsmall}
    j_t = -\,s_t + \phi\,\mathbb{E}_t\!\left[\mathit{SDF}^{\ell}_{t+1}
    \left(q_t\,v_{t+1} + \left(1-q_t\right) j_{t+1}\right)\right].
\end{equation}
Scaling by $A^{\ell}_{t-1}$ exactly as before --- define
$\mathcal{J}^{\ell}_t \equiv A^{\ell}_{t-1} j_t$ and
$S^{\ell}_t \equiv A^{\ell}_{t-1} s_t$ --- and then detrending by $\Lambda_t$
gives:
\begin{equation}\label{eq:J}
    \mathcal{J}_t = -\,S_t + \phi\,\mathbb{E}_t\!\left[\mathit{SDF}_{t+1}\,
    \frac{A_{t-1}}{A_{t}}\,\frac{1}{1+\gamma^{a}}
    \left(q_t\,\mathcal{V}_{t+1} + \left(1-q_t\right)\mathcal{J}_{t+1}\right)\right],
\end{equation}
in the model:
\modelcode{J = -S+phi*(SDF(+1)*A(-1)/A*1/(1+gammaa)*(q*V(+1)+(1-q)*J(+1)));}
The $A$-scaling of $S_t$ is also what squares the market-clearing
adoption cost: total spending is (number of unadopted technologies) $\times\,s_t
= \left(Z_{t-1}/A_{t-1}-1\right) S_t$, the term in the aggregate-demand equation.

\subsection{The FOC for Adoption Expenditure}

The adopter chooses $S_t$ to maximize $\mathcal{J}_t$ in \eqref{eq:J}, taking into
account that spending raises the adoption probability through \eqref{eq:q}. From
\eqref{eq:q},
\begin{equation*}
    \frac{\partial q_t}{\partial S_t} = \varsigma\,\frac{q_t}{S_t}.
\end{equation*}
The FOC is:
\begin{equation*}
    \frac{\partial \mathcal{J}_t}{\partial S_t} = 0
    \;\Leftrightarrow\;
    -1 + \phi\,\mathbb{E}_t\!\left[\mathit{SDF}_{t+1}\,
    \frac{A_{t-1}}{A_{t}}\,\frac{1}{1+\gamma^{a}}\,
    \varsigma\,\frac{q_t}{S_t}
    \left(\mathcal{V}_{t+1} - \mathcal{J}_{t+1}\right)\right] = 0.
\end{equation*}
Multiplying through by $S_t$:
\begin{equation}\label{eq:Sfoc}
    S_t = \varsigma\, q_t\, \phi\,\mathbb{E}_t\!\left[\mathit{SDF}_{t+1}\,
    \frac{1}{1+\gamma^{a}}\,\frac{A_{t-1}}{A_{t}}
    \left(\mathcal{V}_{t+1} - \mathcal{J}_{t+1}\right)\right],
\end{equation}
in the model:
\modelcode{varsigma*q*phi*SDF(+1)/(1+gammaa)*A(-1)/A*(V(+1)-J(+1)) = S;}
\eqref{eq:Sfoc} equates the marginal cost of adoption effort --- one unit of
the final good --- to its marginal value: the probability gain $\varsigma q_t/S_t$
times the discounted adoption surplus $\mathcal{V}_{t+1} - \mathcal{J}_{t+1}$.

\subsection{Aside: Where Each Growth Factor Comes From}

The technology block carries three distinct correction factors, with three
distinct origins:
\begin{itemize}
    \item $A_{t-1}/A_t$ --- a \emph{scaling} conversion, not detrending: the
    values and $S_t$ are per-variety objects scaled by the adopted stock, and the
    stock advances one period inside the continuation value.
    \item $1/(1+\gamma^{a})$ --- the detrending of the technology stocks by their
    \emph{own} trend, which exceeds the output trend $g$ because technology
    enters production only through $A_{t-1}^{\vartheta-1}$ (see the note in
    \texttt{../tech-growth-rate/}).
    \item no $g$ --- the values grow with the output trend $\Lambda_t$, and the
    level SDF equals the detrended one divided by $g_{t+1}$; the two growth
    factors cancel, so the model's $\mathit{SDF}_{t+1}$ appears bare (the same
    cancellation as in the pricing recursions,
    Subsection~\ref{subsec:prg}).
\end{itemize}

%========================================================================
\section{Government: Fiscal and Monetary Policy}\label{sec:gov}
%========================================================================

The government block contains no optimization; it is a set of rules and
identities. This section states them so that the note covers the model in full.

\subsection{Spending Instruments}

Each of the four expenditure instruments is its steady-state level plus an
innovation scaled by steady-state demand; persistence comes from the shock
paths set per experiment, not from the rules themselves:
\modelcode{Gc = Gcy*ydss+ydss*epsi\_gc;\\
Igi = Igiy*ydss+ydss*epsi\_igi;\\
Ige = Igey*ydss+ydss*epsi\_ige;\\
Grd = Grdy*ydss+ydss*epsi\_grd;}
In the budget-neutral reallocation experiments the consumption innovation
offsets the investment ones, \texttt{epsi\_gc} $= -($\texttt{epsi\_igi} $+$
\texttt{epsi\_ige} $+$ \texttt{epsi\_grd}$)$; the debt-financed expansions move a
single instrument with no offset.

\subsection{Public Capital and Efficiency Gaps}

Public investment accumulates into infrastructure and human-capital stocks, with
only the productive fraction $1-e$ of each unit spent arriving in the stock (the
$g$ on the current stock is the stationarization):
\modelcode{Kg*g = (1-delta)*Kg(-1)+(1-eGI)*Igi;\\
Kge*g = (1-delta)*Kge(-1)+(1-eGE)*Ige;}
The efficiency gaps are constants moved only by the experiment shocks; a
positive shock closes the gap:
\modelcode{eGI = eGI\_ss-epsi\_effgi;\\
eGE = eGE\_ss-epsi\_effge;\\
eGRD = eGRD\_ss-epsi\_effgrd;}

\subsection{Budget, Debt, and the Fiscal Rules}

The budget constraint accumulates real debt --- the $/g$ is the stationarization
of the predetermined stock --- and \texttt{by} is the debt-to-GDP ratio:
\modelcode{b = (R(-1)/PI)*b(-1)/g+Gc+Igi+Ige+Grd+T-tauw*w*N-tauc*C;\\
by = b/y;}
Taxes and transfers respond to the lagged deviation of debt-to-GDP from target:
\modelcode{tauc-taucss = rho\_tauc*(tauc(-1)-taucss)\\
\hspace*{2em}+(1-rho\_tauc)*(gamma\_d\_tauc*(by(-1)-byss))+epsi\_tauc;\\
tauw-tauwss = rho\_tauw*(tauw(-1)-tauwss)\\
\hspace*{2em}+(1-rho\_tauw)*(gamma\_d\_tauw*(by(-1)-byss))+epsi\_tauw;\\
T-STEADY\_STATE(T) = rho\_T*(T(-1)-STEADY\_STATE(T))\\
\hspace*{2em}+(1-rho\_T)*(-gamma\_d\_T*(by(-1)-byss)*ydss);}
The minus sign in the transfer rule means transfers fall when debt is above
target. At the calibration, transfers are the only active feedback:
\texttt{gamma\_d\_tauc} $=$ \texttt{gamma\_d\_tauw} $= 0$.

\subsection{Monetary Policy and the Borrowing Rate}

The central bank follows a standard Taylor rule with smoothing, and the
government borrowing rate tracks the policy rate up to an exogenous spread
shock:
\modelcode{Rmp/Rss = (Rmp(-1)/Rss)\^{}rho\_R\\
\hspace*{2em}*((PI/Piss)\^{}gamma\_pi*(yd/ydss)\^{}gamma\_y)\^{}(1-rho\_R)*exp(epsi\_MP);\\
log(R) = rho\_RG*R(-1)+ (1-rho\_RG)*log(Rmp) + epsi\_spread;}
One dormant oddity in the second equation: the lagged term is linear
\texttt{R(-1)} next to logs. It is inert at the calibration
(\texttt{rho\_RG} $=0$, so the equation reduces to
$\log R_t = \log R^{mp}_t + \varepsilon^{spr}_t$), but it would be dimensionally
inconsistent for any \texttt{rho\_RG} $>0$ (flagged in the equation
reconciliation note in \texttt{../../issues/}).

%========================================================================
\section{Market Clearing}\label{sec:mc}
%========================================================================

Aggregate demand sums private consumption and investment, the four public
spending instruments, and the technology-adoption cost of
Section~\ref{sec:tech}; gross output exceeds demand by the price-dispersion
wedge:
\modelcode{yd = C+Ip+Gc+Igi+Ige+Grd+(Z(-1)/A(-1)-1)*S;\\
y = vp*yd;}
The law of motion of the dispersion term \texttt{vp} is derived and boxed in
Subsection~\ref{subsec:vp}.

%========================================================================
\section{Balanced Growth}\label{sec:bg}
%========================================================================

Trend growth is constant absent shocks:
\modelcode{log(g) = (1-rho\_g)*log(g(-1))+rho\_g*(log(gss))+epsi\_g;}
Note the parameterization: the weight on the \emph{lagged} term is
$1-$\texttt{rho\_g}, so \texttt{rho\_g} (calibrated $0.24$) is the speed of
reversion to steady-state growth, not the persistence. No paper experiment
shocks \texttt{epsi\_g}, so $g_t = g$ throughout.

%========================================================================
\section{Auxiliary, Reporting, and Steady-State Pins}\label{sec:aux}
%========================================================================

The remaining equations define reporting variables --- they feed nothing back
into the equilibrium. \texttt{G} is total government spending, \texttt{rreal}
the ex-post real rate, and the last three express fiscal aggregates as shares of
steady-state GDP (a constant denominator, so the panels show pure accumulation):
\modelcode{G = Gc+Igi+Ige+Grd;\\
rreal = R/PI;\\
pdef\_yss = (Gc+Igi+Ige+Grd+T-tauw*w*N-tauc*C)/ydss;\\
T\_yss = T/ydss;\\
by\_yss = b/ydss;}
Finally, five constants are carried as variables pinned to their steady-state
values: the endogenously calibrated scales (\texttt{omega} in labor disutility,
\texttt{chiH} in human-capital accumulation, \texttt{kappaprob} $=q_0$ in the
adoption probability) and the reference levels (\texttt{Rss}, \texttt{ydss})
used inside the policy rules:
\modelcode{omega = STEADY\_STATE(omega);\\
Rss = STEADY\_STATE(R);\\
ydss = STEADY\_STATE(yd);\\
chiH = STEADY\_STATE(chiH);\\
kappaprob = STEADY\_STATE(kappaprob);}
With these, every equation of \texttt{models/model\_block.modpart} has appeared
in a box exactly once.

\end{document}
